% © 2026 Ing. David Jaroš — CC BY-NC-ND 4.0
%
% This work is licensed under a Creative Commons Attribution-NonCommercial-NoDerivatives
% 4.0 International License (CC BY-NC-ND 4.0).
%
% License History: Earlier drafts (up to v0.3) were released under CC BY 4.0.
% From v0.4 onward, all material is released under CC BY-NC-ND 4.0 to protect
% the integrity of the theoretical work during ongoing academic development.
%
% See LICENSE.md for full license text.

% Hecke Bridge: Motivation for the ℂ⊗ℍ ↔ Modular Forms Connection
% Status: [MOTIVATED CONJECTURE — algebraic connection OPEN]
% Date: 2026-03-07

\ifdefined\INCLUDEMODE
\else
  \documentclass[12pt]{article}
  \usepackage[a4paper, margin=2.5cm]{geometry}
  \usepackage{amsmath, amssymb, amsthm}
  \usepackage{hyperref}
  \usepackage{xcolor}
  \usepackage{mdframed}
  \usepackage{tcolorbox}

  \newtheorem{theorem}{Theorem}[section]
  \newtheorem{lemma}[theorem]{Lemma}
  \newtheorem{proposition}[theorem]{Proposition}
  \theoremstyle{definition}
  \newtheorem{definition}[theorem]{Definition}
  \theoremstyle{remark}
  \newtheorem{remark}[theorem]{Remark}

  \newmdenv[backgroundcolor=yellow!20, linecolor=orange, linewidth=1.5pt]{gapbox}
  \newmdenv[backgroundcolor=green!15, linecolor=green!60!black, linewidth=1.5pt]{resultbox}
  \newmdenv[backgroundcolor=blue!10, linecolor=blue!60, linewidth=1.5pt]{motivatebox}

  \tcbuselibrary{skins,breakable}
  \newtcolorbox{statusbox}[1]{
    colback=yellow!8!white, colframe=orange!60!black,
    title={\textbf{#1}}, fonttitle=\bfseries, breakable
  }

  \title{\textbf{Hecke Bridge: Motivation for the $\mathbb{C}\otimes\mathbb{H}
         \leftrightarrow$ Modular Forms Connection}\\
         \large A Motivating Sketch (not a Derivation)}
  \author{David Jaroš}
  \date{2026-03-07}

  \begin{document}
  \maketitle

  \begin{abstract}
  The Hecke conjecture in UBT asserts that the three fermion generations correspond
  to specific modular forms at prime $p = 137$ (and mirror forms at $p = 139$).
  This document presents a \emph{motivating sketch} --- not a rigorous derivation ---
  for why the $\mathbb{C}\otimes\mathbb{H}$ algebraic structure of UBT should select
  modular forms of specific weights $k = 2, 4, 6$ and levels related to
  $\dim_{\mathbb{R}}(\mathbb{C}\otimes\mathbb{H}) = 8$.
  All connections are conjectural; explicit open questions are listed.

  \textbf{Status: [MOTIVATED CONJECTURE --- algebraic connection OPEN]}
  \end{abstract}
\fi

\section{Introduction: The Hecke Conjecture}
\label{sec:intro}

\begin{statusbox}{Status: MOTIVATED CONJECTURE — algebraic connection OPEN}
This document presents motivating arguments, not proofs.  The connection between
$\mathbb{C}\otimes\mathbb{H}$ and modular forms is a conjecture motivated by
numerical observations (unique global hit at $p = 137$ and $p = 139$) and
algebraic suggestiveness.  No rigorous derivation exists.  All open questions
are explicitly listed in Section~\ref{sec:open}.
\end{statusbox}

The Hecke conjecture (see \texttt{reports/hecke\_lepton/}) states that the three
fermion generation mass ratios $m_e : m_\mu : m_\tau$ are encoded in the Hecke
eigenvalues $a_p(f)$ of three specific cuspidal newforms:
\begin{align}
  f_1 &= \text{76.2.a.a} \quad (k=2,\, N=76,\, a_{137} = -11), \\
  f_2 &= \text{7.4.a.a}  \quad (k=4,\, N=7,\,  a_{137} = +2274), \\
  f_3 &= \text{208.6.a.?}\quad (k=6,\, N=208,\, a_{137} = -38286).
\end{align}
The question addressed here is: \emph{why these specific modular forms?}
What property of the $\mathbb{C}\otimes\mathbb{H}$ algebra selects weights
$k = 2, 4, 6$ and levels $N = 76, 7, 208$?

\section{Motivation 1: Weights $k = 2, 4, 6$ from the Action $S[\Theta]$}
\label{sec:weights}

\subsection{Kinetic term is quadratic}

The UBT action is:
\begin{equation}
  S[\Theta] = \int (\partial_\mu\Theta^\dagger)(\partial^\mu\Theta)\sqrt{-g}\,d^4x.
\end{equation}
The kinetic term is \emph{quadratic} in $\Theta$.  In the language of modular forms,
weight $k$ of a cusp form $f$ corresponds to the degree of the associated $L$-function.
The identification:
\begin{itemize}
  \item $k = 2$: kinetic term (quadratic in $\Theta$) $\to$ first generation,
  \item $k = 4$: quartic interaction (four-$\Theta$ vertex) $\to$ second generation,
  \item $k = 6$: sextic term (six-$\Theta$ vertex) $\to$ third generation,
\end{itemize}
suggests that the $n$-th generation corresponds to the $n$-th power of $\Theta$ in the
action, mapped to modular weight $k = 2n$.

\begin{motivatebox}
\textbf{Motivation}: Weight $k = 2n$ corresponds to the $n$-th $\psi$-mode of $\Theta$.
The kinetic ($k=2$), quartic ($k=4$), and sextic ($k=6$) terms in $S[\Theta]$ naturally
select forms of weight 2, 4, 6 as the relevant modular objects associated with the
three fermion generations.
\end{motivatebox}

This is a \emph{motivating analogy}, not a theorem.  The precise relation between
powers of $\Theta$ in the action and modular weights requires establishing a
dictionary between the $\psi$-mode expansion of $\Theta$ and the Hecke algebra action.

\section{Motivation 2: Level $N = 7$ and $\dim_{\mathbb{R}}(\mathbb{C}\otimes\mathbb{H})$}
\label{sec:level}

\subsection{The suggestive coincidence}

The modular curve $\Gamma_0(7)$ has index:
\begin{equation}
  \mu(\Gamma_0(7)) = 7\cdot\prod_{p|7}\left(1 + \frac{1}{p}\right) = 7\cdot\frac{8}{7} = 8.
\end{equation}
This equals $\dim_{\mathbb{R}}(\mathbb{C}\otimes\mathbb{H}) = 8$.

Additionally, $N = 7$ is the smallest prime for which $S_4(\Gamma_0(N))$ (the space
of cusp forms of weight 4) is non-trivial:
\begin{equation}
  \dim S_4(\Gamma_0(7)) = 1,
\end{equation}
giving a unique newform $f_2 = \text{7.4.a.a}$ at $(N,k) = (7,4)$.

\begin{motivatebox}
\textbf{Motivation}: $N = 7$ is selected for the second generation ($k = 4$)
because it is the smallest level with a non-trivial $k=4$ cusp form, and its
modular index equals $\dim_{\mathbb{R}}(\mathbb{C}\otimes\mathbb{H}) = 8$.
This suggests the biquaternionic dimension $8$ controls the minimal level.
\end{motivatebox}

The connection to the Klein quartic (genus-3 curve $x^3y + y^3z + z^3x = 0$
with automorphism group $\mathrm{PSL}(2,\mathbb{F}_7)$ of order $168 = 8 \times 21$)
may be relevant: the genus of $X_0(7)$ is 0, but higher-genus covers of
$\mathcal{H}/\Gamma_0(7)$ relate to $\mathrm{PSL}(2,\mathbb{F}_7)$.

\subsection{Why Level $N = 7$ for the second generation?}

This is an \textbf{open question}.  The current answer is: ``$N = 7$ is the minimal
level with a unique non-trivial $k=4$ newform; the modular index matches the
biquaternionic dimension; this is suggestive but not proved.''

\section{Motivation 3: Trivial Character $\chi = 1$ from $\mathbb{Z}$-Rationality}
\label{sec:character}

All six Hecke forms (Set A and Set B) have trivial character $\chi = 1$, meaning
their Fourier coefficients $a_n$ are integers.

The $\mathbb{C}\otimes\mathbb{H}$ algebra is defined over $\mathbb{Z}$
(as a Clifford algebra: generators $e_1, e_2, e_3, e_4$ with integer structure
constants).  The natural modular forms associated to such an algebra should have
integer Fourier coefficients --- i.e., trivial character.

\begin{motivatebox}
\textbf{Motivation}: $\mathbb{C}\otimes\mathbb{H}$ is defined over $\mathbb{Z}$
$\Rightarrow$ UBT naturally selects $\mathbb{Q}$-rational modular forms
(trivial character, integer Hecke eigenvalues).
All six identified Hecke forms satisfy this constraint.
\end{motivatebox}

\section{Motivation 4: Three Forms (Not One, Not Five)}
\label{sec:three}

Why exactly three modular forms (one per generation)?

The number of independent fermion generations in UBT is explained by three independent
$\psi$-modes $\Theta_0, \Theta_1, \Theta_2$ (see \texttt{research\_tracks/three\_generations/}).
Each mode has its own characteristic energy scale and interacts with its own
``generation-specific'' Hecke form.

The triple structure $(k=2, k=4, k=6)$ corresponds to:
\begin{itemize}
  \item $\Theta_0$ (mode 0, lightest = electron family): $k = 2$,
  \item $\Theta_1$ (mode 1, muon family): $k = 4$,
  \item $\Theta_2$ (mode 2, tau family): $k = 6$.
\end{itemize}
The weights $k = 2, 4, 6$ form an arithmetic sequence with step 2, consistent
with the Kaluza--Klein mode expansion in the $\psi$-direction.

\section{Open Questions}
\label{sec:open}

The following are the most important open questions for the Hecke bridge:

\begin{enumerate}
  \item \textbf{How to derive $N = 7$ from $\mathbb{C}\otimes\mathbb{H}$ geometry?}
    The index $\mu(\Gamma_0(7)) = 8 = \dim_{\mathbb{R}}(\mathbb{C}\otimes\mathbb{H})$
    is suggestive but not proved.  A rigorous connection would require identifying
    $\Gamma_0(7)$ as the modular symmetry group of the $\psi$-bundle over some
    arithmetic surface associated to $\mathbb{C}\otimes\mathbb{H}$.

  \item \textbf{Why exactly three forms (not one, not five)?}
    The triple structure is motivated by three $\psi$-modes, but a derivation
    from first principles (without appealing to the observed number of generations)
    is lacking.

  \item \textbf{Shimura lifts between $k = 2$, $k = 4$, $k = 6$ forms}:
    Are there Shimura lift maps connecting the three Set A forms?
    Such lifts would provide an algebraic structure linking the three generations.

  \item \textbf{Algebraic derivation of the specific forms}:
    The forms $76.2.a.a$, $7.4.a.a$, $208.6.a.?$ were found by numerical search.
    A derivation from $\mathbb{C}\otimes\mathbb{H}$ algebra would require identifying
    the $L$-function of $\mathbb{C}\otimes\mathbb{H}$ over $\mathbb{Z}$ and matching
    it to the product $L(f_1, s) \cdot L(f_2, s) \cdot L(f_3, s)$.

  \item \textbf{Mirror sector levels ($N = 195$, $N = 50$, $N = 54$)}:
    The Set B levels do not share the $\mu(\Gamma_0(N)) = 8$ property.
    Why are these specific levels selected for the mirror sector?

  \item \textbf{Connection to Klein quartic and $\mathrm{PSL}(2, \mathbb{F}_7)$}:
    The automorphism group of the Klein quartic has order $168 = 8 \times 21$.
    Is there a deeper connection between the Klein quartic and the $k=4$, $N=7$ form?
\end{enumerate}

\section{Summary}

\begin{center}
\renewcommand{\arraystretch}{1.3}
\begin{tabular}{llp{5cm}}
\toprule
\textbf{Structural element} & \textbf{Motivation} & \textbf{Status} \\
\midrule
Weights $k = 2, 4, 6$ & $n$-th $\psi$-mode $\leftrightarrow$ weight $k=2n$ &
  [MOTIVATED CONJECTURE] \\
Level $N = 7$, $\mu(\Gamma_0(7)) = 8$ & $\dim_{\mathbb{R}}(\mathbb{C}\otimes\mathbb{H})=8$ &
  [MOTIVATED CONJECTURE] \\
Trivial character $\chi = 1$ & $\mathbb{C}\otimes\mathbb{H}$ over $\mathbb{Z}$ &
  [MOTIVATED CONJECTURE] \\
Three forms (not one, not five) & Three $\psi$-modes &
  [MOTIVATED CONJECTURE] \\
Derivation of specific forms & — &
  \textbf{OPEN} \\
Derivation of specific levels & — &
  \textbf{OPEN} \\
\bottomrule
\end{tabular}
\end{center}

\noindent\textbf{Date}: 2026-03-07 \\
\textbf{Author}: David Jaroš

\begin{thebibliography}{9}
\bibitem{HeckeLMFDB}
The LMFDB Collaboration, \textit{The L-functions and Modular Forms Database},
\texttt{https://www.lmfdb.org}, 2024.

\bibitem{HeckeLepton}
D.\ Jaroš, \textit{Hecke Conjecture: Lepton Mass Ratios from Modular Forms},
\texttt{reports/hecke\_lepton/}, 2026.

\bibitem{Shimura}
G.\ Shimura, \textit{Introduction to the Arithmetic Theory of Automorphic Functions},
Princeton University Press, 1971.
\end{thebibliography}

\ifdefined\INCLUDEMODE
\else
\end{document}
\fi
